\section{Koide Phase Overlap and Finite-Gap Constraints}
\label{sec:v4110-koide-overlap}

The original OP~3 asked for the phase of the charged-lepton mass
spectrum. This section examines a conditional readout of several
phase-shifted preparations in a common generation basis in frame M.
Each preparation has nonnegative roots, a common mean \(a>0\), and
the exact Koide ratio \(2/3\). Let
\[
 x_{jk}=a\bigl[1+\sqrt2\,\operatorname{Re}(\omega^k e^{i\delta_j})\bigr],
 \qquad \omega=e^{2\pi i/3},\qquad k=0,1,2.
\]
If the readout averages the roots with nonnegative weights summing to
one before squaring, it is determined by \(Z=\sum_j w_j e^{i\delta_j}\).
In particular, the resulting Koide ratio is \((1+|Z|^2)/3\), and the
reduced phase is \(\arg Z\) for \(Z\ne0\). A pair with equal weights
at \(\delta_0\pm h\), for \(|h|<\pi/2\), retains the central phase
and reduces the radius by \(\cos h\). Averaging the squared roots
instead adds the weighted variance of each root to the corresponding
mass. This comparison makes the order of the readout operations
explicit.

\subsection{Exact poles and the positive-root arc}

Consider a cyclic comparison with static root eigenvalues \(x_k\)
of mean one, direct diagonal entry \(a_0=1+\sqrt2\), and a common
positive intermediate entry \(D\). The lower eigenvalue \(\ell_k\)
of the block with off-diagonal squared magnitude \(D(a_0-x_k)\)
obeys
\[
 x_k=\ell_k+\frac{\ell_k(a_0-\ell_k)}{D}.
\]
This relation allows the static target to be reconstructed from the
mass ratios at a prescribed \(D\). In the fixed-label chamber about the
positive-radius reference direction the admissible reduced phases lie
in \([-\pi/12,\pi/12]\). The reconstructed complex target must lie in
the convex hull of this arc. With \(N\) equal weights and fixed mean sine,
the smallest mean cosine is attained with all but at most one of the
sines at the endpoints. This follows from concavity of
\(y\mapsto\sqrt{1-y^2}\) on the sine interval. Jensen's inequality gives
the upper bound. Continuity on the connected set of admissible sine
tuples gives every intermediate value.

Using the reference masses \(0.51099895\), \(105.6583755\), and
\(1776.93\) MeV, the comparison at \(D=68\) admits only \(N=15\) among
\(2,\ldots,16\). At \(D=100\) it admits \(N=10,\ldots,16\).
Repeated phases are allowed. The masses, the chosen intermediate
entry and the assumed readout specify this inverse calculation.

\subsection{Conditions for a transport interpretation}

A relative transport phase could enter an effective coupling if two
paths contribute to the same matrix element before pole determination.
A phase change confined to the prepared probe changes its spectral
weights while leaving the operator poles fixed. Likewise, a common
factor multiplying the entire mass-root operator changes the scale
while leaving the mass ratios fixed. The native generation operator,
the coherent end-channel map, the path selection and the intermediate
scale remain open. The structural shell-transition question bearing the
later OP~3 label remains distinct from this spectral question.
